% !TEX encoding = UTF-8 Unicode
\documentclass[11pt,a4paper]{article}

\usepackage[utf8]{inputenc}
\usepackage[T1]{fontenc}
\usepackage{geometry}
\usepackage{hyperref}
\usepackage{enumitem}
\usepackage{booktabs}
\usepackage{array}
\usepackage{tabularx}
\usepackage{xcolor}
\usepackage{colortbl}
\usepackage{fancyhdr}
\usepackage{titlesec}
\usepackage{parskip}
\usepackage{ragged2e}
\usepackage{tcolorbox}
\tcbuselibrary{breakable, skins}

\geometry{a4paper, top=2.2cm, bottom=2.2cm, left=2.5cm, right=2.5cm}

\definecolor{accentblue}{RGB}{30, 80, 160}
\definecolor{warnbg}{RGB}{255, 248, 235}
\definecolor{warnborder}{RGB}{190, 80, 0}
\definecolor{tipbg}{RGB}{240, 248, 240}
\definecolor{tipborder}{RGB}{50, 130, 60}
\definecolor{darkgray}{RGB}{70, 70, 70}
\definecolor{medgray}{RGB}{130, 130, 130}
\definecolor{taskbg}{RGB}{242, 244, 248}
\definecolor{taskborder}{RGB}{30, 80, 160}
\definecolor{labelcolor}{RGB}{40, 40, 100}

\pagestyle{fancy}
\fancyhf{}
\fancyhead[L]{\small\textcolor{medgray}{Feladatlap -- AI képzés}}
\fancyhead[R]{\small\textcolor{medgray}{Sapientia EMTE -- dr.\ Pál László}}
\fancyfoot[C]{\small\textcolor{medgray}{\thepage}}
\renewcommand{\headrulewidth}{0.3pt}

\titleformat{\section}{\large\bfseries\color{accentblue}}{\thesection}{0.6em}{}%
  [\vspace{-2pt}\textcolor{accentblue}{\rule{\linewidth}{0.6pt}}]
\setcounter{secnumdepth}{0}
\titleformat{\subsection}{\normalsize\bfseries\color{darkgray}}{}{0em}{}
\titlespacing{\section}{0pt}{18pt}{8pt}
\titlespacing{\subsection}{0pt}{12pt}{4pt}

\tcbset{
  taskstyle/.style={
    breakable, enhanced,
    colback=taskbg, colframe=taskborder,
    boxrule=1pt, arc=3pt,
    left=10pt, right=10pt, top=10pt, bottom=8pt,
    attach boxed title to top left={xshift=10pt, yshift=-\tcboxedtitleheight/2},
    boxed title style={
      colback=taskborder, colframe=taskborder,
      arc=2pt, boxrule=0pt,
      left=6pt, right=6pt, top=2pt, bottom=2pt
    },
    coltitle=white, fonttitle=\footnotesize\bfseries\sffamily,
    before upper=\RaggedRight,
  },
  warnstyle/.style={
    breakable, colback=warnbg, colframe=warnborder,
    boxrule=0.5pt, leftrule=3pt, arc=2pt,
    left=8pt, right=8pt, top=5pt, bottom=5pt,
    fontupper=\small\sffamily,
  },
  tipstyle/.style={
    breakable, colback=tipbg, colframe=tipborder,
    boxrule=0.5pt, leftrule=3pt, arc=2pt,
    left=8pt, right=8pt, top=5pt, bottom=5pt,
    fontupper=\small\sffamily,
  },
}

\newcommand{\taskbox}[2]{%
  \begin{tcolorbox}[taskstyle, title=#1]
  #2
  \end{tcolorbox}\vspace{4pt}}

\newcommand{\warn}[1]{%
  \vspace{3pt}%
  \begin{tcolorbox}[warnstyle]%
  \textbf{\textcolor{warnborder}{Figyelem:}} #1%
  \end{tcolorbox}\vspace{2pt}}

\newcommand{\tip}[1]{%
  \vspace{3pt}%
  \begin{tcolorbox}[tipstyle]%
  \textbf{\textcolor{tipborder}{Tipp:}} #1%
  \end{tcolorbox}\vspace{2pt}}

\hypersetup{colorlinks=true, linkcolor=accentblue, urlcolor=accentblue}

\begin{document}

\thispagestyle{empty}
\begin{center}
  {\LARGE\bfseries Generatív AI és alkalmazások}\\[5pt]
  {\normalsize Feladatlap}\\[1pt]
  {\small\color{medgray} dr.\ Pál László, Sapientia EMTE, 2026.\ március 6--7.}
\end{center}

%=============================================================
\section{Előkészületek}
%=============================================================

A képzés során két fő eszközt használunk:

\begin{enumerate}[leftmargin=1.5em]
  \item \textbf{Google Gemini} -- \url{https://gemini.google.com}\\
    Jelentkezzetek be a saját e-mail címetekkel.
  \item \textbf{Google Colab} -- \url{https://colab.research.google.com}\\
    Szintén a saját e-mail címetekkel. Erre a 2.\ és 3.\ feladatnál lesz szükség.
\end{enumerate}

\noindent A feladatok anyagai elérhetők:
\begin{itemize}[leftmargin=1.5em]
  \item \textbf{GitHub:} \url{https://github.com/pallaszlo/ai-kepzes}
  \item \textbf{Google Classroom:} \url{https://classroom.google.com/u/0/c/ODQ1Njg1MjM4MDIw}\\
    Kód: \texttt{lgwr5qh3}
\end{itemize}

\noindent Innen tudjátok letölteni az összes szükséges fájlt (feladatleírások, notebookok, adatfájlok).

%=============================================================
\section{A képzés programja}
%=============================================================

\definecolor{headerbg}{RGB}{30, 80, 160}
\definecolor{headertext}{RGB}{255, 255, 255}
\definecolor{rowalt}{RGB}{245, 247, 252}
\definecolor{breakrow}{RGB}{255, 250, 240}

\begingroup
\renewcommand{\arraystretch}{1.35}
\begin{center}
\begin{tabularx}{\textwidth}{c X c l}
  \toprule
  \rowcolor{headerbg}
  \textcolor{headertext}{\textbf{\#}} &
  \textcolor{headertext}{\textbf{Tevékenység}} &
  \textcolor{headertext}{\textbf{Időtartam}} &
  \textcolor{headertext}{\textbf{Eszköz}} \\
  \midrule
  0. & Bevezető előadás & $\sim$30 perc & Prezentáció \\
  \rowcolor{rowalt}
  1. & Promptolás (Prompt engineering) & $\sim$70 perc & Gemini \\
  \rowcolor{breakrow}
     & \textit{--- Szünet ---} & \textit{10 perc} & \\
  2. & AI alapú részvényelemzés & $\sim$100 perc & Colab + Gemini \\
  \rowcolor{breakrow}
     & \textit{--- Ebédszünet ---} & \textit{30 perc} & \\
  \rowcolor{rowalt}
  3. & Airbnb szálláshelyek elemzése & $\sim$115 perc & Colab + Gemini \\
  \midrule
  \rowcolor{headerbg}
     & \textcolor{headertext}{\textbf{Összesen}} & \textcolor{headertext}{\textbf{$\sim$355 perc}} & \\
  \bottomrule
\end{tabularx}
\end{center}
\endgroup

%=============================================================
\section{1.\ feladat -- Promptolás (Prompt engineering) \hfill {\normalsize\color{medgray}$\sim$70 perc}}
%=============================================================

\taskbox{Teendők}{
\begin{enumerate}[leftmargin=1.5em]
  \item Nyissátok meg a \textbf{Gemini}-t a böngészőben: \url{https://gemini.google.com}
  \item Töltsétek le a \texttt{02\_promptolas} mappából a \textbf{Prompt engineering a gyakorlatban.pdf} fájlt.
  \item Nyissátok meg a PDF-et és kövessétek az utasításokat lépésről lépésre.
  \item A feladatleírásban megadott promptokat próbáljátok ki a Gemini felületen.
\end{enumerate}
}


%=============================================================
\section{2.\ feladat -- Részvényelemzés \hfill {\normalsize\color{medgray}$\sim$100 perc}}
%=============================================================

\taskbox{Teendők}{
\begin{enumerate}[leftmargin=1.5em]
  \item Nyissátok meg a \textbf{Google Colab}-ot: \url{https://colab.research.google.com}\\
    Jelentkezzetek be a saját e-mail címetekkel.
  \item Nyissátok meg az előkészített notebookot (\textbf{AI\_alapu\_reszvenyelemzes.ipynb}) a Colab-ban.
  \item Töltsétek le a feladatleírást: \textbf{AI\_alapu\_reszvenyelemzes.pdf} -- és kövessétek az instrukciókat.
  \item A futtatáshoz három adatfájl szükséges (a \texttt{03\_reszveny\_elemzes/adatok} mappában):
    \begin{itemize}
      \item \texttt{AAPL\_adatok.csv}
      \item \texttt{AAPL\_hirek.csv}
      \item \texttt{GSPC\_adatok.csv}
    \end{itemize}
  \item A notebookban a kód megpróbálja letölteni az adatokat a Yahoo Finance-ról. Ha ez nem sikerül, töltsétek fel manuálisan a fenti CSV fájlokat a Colab-ba.
\end{enumerate}
}

%=============================================================
\section{3.\ feladat -- Airbnb elemzés \hfill {\normalsize\color{medgray}$\sim$115 perc}}
%=============================================================

\taskbox{Teendők}{
\begin{enumerate}[leftmargin=1.5em]
  \item Nyissátok meg a \textbf{Google Colab}-ot (ha még nem lenne megnyitva).
  \item Nyissátok meg az előkészített notebookot (\textbf{airbnb\_elemzes.ipynb}) a Colab-ban.
  \item Töltsétek le a feladatleírást: \textbf{Airbnb\_elemzes.pdf} -- és kövessétek az instrukciókat.
  \item Az elemzéshez szükséges adatfájl (a \texttt{04\_airbnb\_elemzes/adatok} mappában):
    \begin{itemize}
      \item \texttt{airbnb\_europe.csv}
    \end{itemize}
  \item Töltsétek fel az adatfájlt a Colab-ba, majd kövessétek a notebook és a PDF lépéseit.
\end{enumerate}
}

\vfill
\begin{center}
  \small\color{medgray} Jó munkát!
\end{center}

\end{document}
